\section*{Discussion}

GraphPop demonstrates that representing population genomic data as a graph---where genotypes, annotations, and computed statistics coexist as a single queryable structure---does not merely accelerate existing analyses but opens investigative pathways that produce new biological insight. The rice 3K application illustrates this concretely: the first systematic genome-wide quantification of $\pi_N/\pi_S$ across all 12 cultivated rice subpopulations, and the first genome-wide detection of adaptive protein divergence between ecotypes, were both enabled by graph-native annotation conditioning. These are not demonstrations of capability; they are biological findings about the cost of domestication and adaptive divergence that were inaccessible to the flat-file tool ecosystem---not because the underlying statistics are novel, but because no existing tool integrates statistics with functional annotations at the scale required.

\textbf{Three architectural properties underlie these capabilities.} First, graph-native computation replaces matrix operations with pre-aggregated node properties, yielding $O(V \times K)$ FAST PATH performance and 87\% memory reduction for FULL PATH analyses. The 12--1{,}614$\times$ speedups over scikit-allel are not incremental: they reflect a fundamentally different computational model (reading structured node properties vs.\ parsing genotype matrices). Second, persistent queryable results transform the database into a cumulative analytical record. The multi-statistic convergence analysis at rice domestication loci---querying already-computed iHS, XP-EHH, PBS, and Fay \& Wu's $H$ on the same nodes---would require custom file-merging scripts in any classical workflow. In GraphPop, it is a single Cypher query. Third, annotation-conditioned queries exploit graph topology to filter variants by consequence or pathway before statistical computation, making analyses like $\pi_N/\pi_S$ and missense/synonymous Fst a one-parameter change rather than a multi-step pipeline.

\textbf{Comparison with Hail.} Hail\cite{hail2024}, built on Apache Spark, is the most comparable large-scale platform. Hail's matrix table model achieves scalability through distributed computing but does not natively represent annotation topology: conditioning on consequence type requires explicit filter--join operations, not graph traversal. Hail does not persist computed statistics as properties on the same data objects---each analysis produces a separate table. GraphPop trades horizontal scalability for a richer data model where the analytical record is integrated with the data and annotations it derives from.

\textbf{Publicly available graph databases.} We deposit pre-built Neo4j graph databases for both the rice 3K (29.6M variants, all computed statistics) and 1000 Genomes (full genome, all autosomes) datasets. These databases contain not only the raw data but all computed results as persistent node properties. Any researcher can load these graphs into a Neo4j instance and immediately query the full analytical record---retrieving, composing, and extending prior results without re-computation. This represents a new mode of data sharing in population genomics: not flat files that must be re-analyzed, but pre-computed analytical records that are immediately queryable.

\textbf{AI agent integration.} GraphPop provides an MCP (Model Context Protocol) server exposing all 12 procedures as tool calls for large language model agents, registered with ToolUniverse (Zitnik Lab, Harvard). This enables AI-driven scientific workflows where an agent can compute statistics, condition on annotations, query convergent signals, and traverse to genes autonomously. The graph database's Cypher interface is particularly well-suited for LLM interaction: queries are declarative, self-documenting, and return structured results. We envision a future where AI agents routinely interrogate graph-native population genomic databases to formulate and test evolutionary hypotheses.

\textbf{Limitations.} GraphPop requires a Neo4j deployment, adding infrastructure complexity. Pairwise LD is slower than PLINK~2 with native binary files. PCA, kinship/IBS/IBD, demographic inference, and simulation interfaces are not yet implemented. The import pipeline requires pre-called VCF and VEP annotations; joint calling must be performed upstream.

\textbf{Future directions.} Planned extensions include graph-native PCA and kinship matrices, integration with demographic inference tools via persistent joint SFS, forward simulation interfaces, and longitudinal cohort tracking. GraphMana, a companion data management platform (separate publication), provides multi-format export, cohort management, and quality control.

GraphPop shows that the same graph-native approach that has transformed knowledge graphs, social networks, and biological pathway databases can fundamentally reshape population genomics---not just as a faster engine, but as a new investigative lens that reveals biology invisible to tools operating on flat files.
